\scanpage{16}
\begin{center}
COURBES ELLIPTIQUES : FORMULAIRE

d'après J. Tate\\
mis au goût du jour par P. Deligne
\end{center}

1.\qquad \uline{Formules générales}.

Dans ce formulaire, on appellera courbe de genre $1$ sur un schéma de base $\mathrm{S}$ un morphisme propre et plat de présentation finie
\[
\mathrm{p}:\mathrm{E}\handwritten{\longrightarrow}\mathrm{S}\qquad,
\]
muni d'une section $\mathrm{e}$ contenue dans l'ouvert de lissité, dont les fibres géométriques sont des courbes réduites irréductibles de \handwritten{\uline}{genre arithmétique}\reversemarginpar\margin{\handwritten{?}}\normalmarginpar $1$. Les fibres géométriques de $\mathrm{p}$ sont donc soit :

a) une courbe elliptique, i.e. une courbe propre, lisse et connexe de genre $1$,

b) une droite projective dont on a identifié deux points distincts (cubique dans $\mathrm{P}_{2}$ à point double ordinaire),

c) une droite projective dont on a identifié deux points infiniment voisins (cubique dans $\mathrm{P}_{2}$ à singularité cuspidale).

On dira dans le cas a) (resp. b), c)) que la fibre est de type elliptique (resp. multiplicatif, resp. additif).

Si $\mathrm{E}$ est une courbe de genre $1$ sur $\mathrm{S}$, on définit un faisceau inversible $\omega$ sur $\mathrm{S}$ par la formule
\[
\omega=\mathrm{e}^{*}\,\Omega^{1}_{\mathrm{E}/\mathrm{S}}\qquad.
\]

La section $\mathrm{e}$ est un diviseur de Cartier relatif de $\mathrm{E}$ sur $\mathrm{S}$ ;
% archive sheet number: 15
\clearpage
